\begin{abstract}
We derive and analyze ordinary differential equations (ODEs) governing
series RLC circuits to show how resistance controls damping and stability.
First-order ODEs governing simplified RC and RL sub-circuits, as well as the
second-order ODE describing the full RLC system, are derived and solved
analytically for all three configurations, and the role of
resistance in controlling circuit damping is examined by varying component
parameters.
The ideal model is then extended numerically to include current-dependent
resistor heating, inductor saturation, and resonant AC forcing.
These nonlinear extensions are phenomenological models chosen to illustrate
how stress-dependent component behavior can alter transient damping.
Our central research question asks: \emph{How do resistors influence the
damping and stability of series RLC circuits as described by ODEs, and what
are the practical implications for engineering applications such as filters
and oscillators?}
The analysis identifies resistance as the parameter that governs whether the
circuit is overdamped, critically damped, or underdamped, while the numerical
results show how non-ideal components can suppress, shift, or amplify
transient oscillations.
\end{abstract}


\begin{IEEEkeywords}
series RLC circuit, ordinary differential equations, damping,
Kirchhoff's Voltage Law, stability analysis, nonlinear simulation,
signal processing.
\end{IEEEkeywords}
